\answer{OR07}
先设 $0<p<1$，记 $q=1-p$。以每次出发时所在地点的雨伞数 $X_t$ 为状态，状态空间为 $\{0,1,\ldots,n\}$。由于下一次从另一地点出发，转移为
\[
 P(0,n)=1,\qquad
 P(i,n-i)=q,\quad P(i,n-i+1)=p\quad(i\ge1).
\]
这条有限链不可约：从 $i<n$ 经一次不带伞、一次带伞可到 $i+1$；从 $i\ge2$ 经一次带伞、一次不带伞可到 $i-1$，而从1还可经 $n$ 到0。这些路径的概率均为正。

候选平稳分布为
\[
 \pi_0=\frac q{n+q},\qquad
 \pi_i=\frac1{n+q}\quad(1\le i\le n).
\]
流入0的概率为 $q\pi_n=\pi_0$。对 $1\le j<n$，流入概率为 $q\pi_{n-j}+p\pi_{n-j+1}=\pi_j$；流入 $n$ 的概率为 $\pi_0+p\pi_1=\pi_n$。这些概率和为1，故确为唯一平稳分布。

只有处于状态0且当次下雨时才会淋雨。有限不可约链的长期平均奖励等于平稳平均奖励，因此所求频率几乎必然为
\[
 \boxed{p\pi_0=\frac{p(1-p)}{n+1-p}}.
\]
该极限与初始的 $k$ 无关。$p=0$ 时从不淋雨；$p=1$ 时至多在第一次出行淋雨，之后每次出发地均有伞，长期频率也为0。上式在这两个端点仍然成立。
